\clearpage
\SourceLeaf{065}\pdfbookmark[1]{Section II.}{sectio-secunda}
\begin{center}
{\large SECTION II.}\par\medskip
\emph{Relations concerning a single position in space.}\par\bigskip
\hypertarget{native-065-article047}{}\pdfbookmark[2]{47.}{article-047}47.
\end{center}
In Section I we dealt with the motion of celestial bodies in their orbits without considering how those orbits are situated in space. To determine this orientation so that we can relate the positions of a celestial body to any other points in space, we evidently need both the orientation of the plane in which the orbit lies relative to a known plane (for example, the plane of the Earth’s orbit, the \emph{ecliptic}) and the position of the apsides in that plane. This is most conveniently treated by spherical trigonometry. We imagine a spherical surface of arbitrary radius with the Sun at its center. On it any plane passing through the Sun traces a great circle, and any straight line drawn from the Sun marks a point. For planes or straight lines that do not pass through the Sun itself, we draw parallel planes and lines through the Sun and regard the corresponding great circles and points on the sphere as representing the original planes and lines as well. We may also suppose the sphere to have what is called an infinitely large radius, in which case parallel planes and lines are represented in the same way.

Unless the plane of the orbit coincides with the plane of the ecliptic, the great circles corresponding to those planes (which we shall also simply call the orbit and the ecliptic) intersect at two points called \emph{nodes}. At one node the body, as seen from the Sun, passes from the southern region across the ecliptic into the northern region. At the other it returns from the northern to the southern region. The former is called the \emph{ascending node}, and the latter the \emph{descending node}. We specify the positions of the nodes on the ecliptic by their distance from the mean
\SourceLeaf{066}vernal equinox (their \emph{longitude}), reckoned in the order of the zodiacal signs. In Fig.~1, let $\AscendingNode$ be the ascending node, $A\AscendingNode B$ a part of the ecliptic, and $C\AscendingNode D$ a part of the orbit. Let the Earth and the celestial body move in the directions from $A$ toward $B$ and from $C$ toward $D$. It is clear that the spherical angle that $\AscendingNode D$ makes with $\AscendingNode B$ can increase from $0$ to $180^0$. Beyond this, $\AscendingNode$ would cease to be the ascending node. We call this angle the \emph{inclination of the orbit} to the ecliptic. Once the orientation of the orbital plane is determined by the longitude of the node and the inclination of the orbit, all that is still needed is the distance of perihelion from the ascending node. We reckon it in the direction of motion itself, and therefore take it to be negative or between $180^0$ and $360^0$ whenever perihelion lies south of the ecliptic. We also note the following expressions. The longitude of any point on the circle of the orbit is reckoned from the point that lies as far back from the ascending node along the orbit as the vernal equinox lies back from that same node along the ecliptic. Hence the \emph{longitude of perihelion} is the sum of the longitude of the node and the distance of perihelion from the node. The \emph{true longitude} of the body \emph{in the orbit} is the sum of the true anomaly and the longitude of perihelion. Finally, the sum of the mean anomaly and the longitude of perihelion is called the \emph{mean longitude}. This last expression can evidently apply only to elliptical orbits.

